\subsection{Résultats}

\subsubsection{Temps de propagation}

Nous mesurons donc 2 signaux sur l'oscilloscope (visible sur \ref{fig:r0} : l'impulsion du generateur (CH1) et l'impulsion passant par le cable coaxial de 100m (CH2).

\begin{figure}[ht]
	\centering
	\includegraphics[scale=0.6]{electrique/r5.png}
	\label{fig:r0}
	\caption{Mesure des impulsions de largeur (width) $w$ = 210 $\mu$$s$ avec R = 0$\ohm$}
\end{figure}

Le delai entre les deux impulsions est definit tel que $\delta$$t$ = $|$$t_1$ - $t_2$$|$ est approximativement, selon Fig. \ref{fig:r0} de : 

\[
        \mathrm{\delta t} \approx 0,5 \mu s 
\]


\subsubsection{Impedence characteristique}

Nous avons mesure grace au pont RLC les valeurs suivantes pour la capacite lineique $c*$ et $l*$ :

\[
	\mathrm{c*} = 10,1 nF \qquad
	\mathrm{l*} = 31,0 \mu H
\]

Ainsi comme nous connaissons la relation liant c* et l* avec l'impedance characteristique :
\[
	\mathrm{Z_c} = \sqrt{\frac{l*}{c*}} = \sqrt{\frac{31,0 \cdot 10^-6}{10,1 \cdot 10^-9}} = 55,5 \pm 0,3
\]

Pour les incertitudes :
\[
\Delta Z_c = \frac{1}{2} \cdot Z_c \sqrt{\left(\frac{\Delta L^*}{L^*}\right)^2 + \left(\frac{\Delta C^*}{C^*}\right)^2}
\]

Donc :
\[
\Delta Z_c = 55,5 \cdot \frac{1}{2} \sqrt{\left(\frac{0,1}{31,0}\right)^2 + \left(\frac{0,1}{10,1}\right)^2} \approx 0,3~\Omega
\]

\subsubsection{Coefficients de reflexion et de transmission}

%Prenons le cas ou nous avons un coefficient de reflexion relativement faible. Nous devons obtenir un comportement tel que sur Fig. \ref{concept.png} tel sont liées les amplitudes et le coefficient de reflexion en Eq. \ref{eq:a}. Considerons les phases 1, 2 et 3. (1) Onde progressive et $A_1$ positive. (2) Onde retrograde avec$A_2$ negative et $|A_2| < |A_1|$. (3) Onde progressive avec $A_3$ positive. 
%\begin{figure}[ht]
%	\centering
%	\includegraphics[scale=0.3]{concept.png}
%	\caption{Evolution d'une impulsion avec faible $\rho$}
%	\label{concept.png}
%\end{figure}
Le chemin sur CH1 est tellement court que nous pouvons observer sur chaque figure une seconde onde $\epsilon_2$ reflexion de $\epsilon_1$. Sur Figs. \ref{fig:r25} et \ref{fig:r50}, $\epsilon_2$ a le meme pattern que $\epsilon_1$, mais avec une amplitude plus basse. Cette amplitude baisse sur Fig. \ref{fig:r50} en accord avec la relation Eq. \ref{eq:a}.

\begin{figure}[h!]
    \centering
    % Première figure
    \begin{minipage}[b]{0.48\linewidth}
        \centering
        \includegraphics[width=\linewidth]{electrique/r25.png}
            \caption{$Z_L$ = 25 $\ohm$}
        \label{fig:r25}
    \end{minipage}
    \hfill
    % Deuxième figure
    \begin{minipage}[b]{0.48\linewidth}
        \centering
        \includegraphics[width=\linewidth]{electrique/christophe/r_50.png}
	    \caption{$Z_L$ = 50 $\ohm$ Ref \cite{christophe}}
        \label{fig:r50}
    \end{minipage}
\end{figure}

En outre, lorsque $Z_L$ devient plus grand que $Z_C$ (voir Fig. \ref{fig:r100}), nous observons l'onde reflechie avec $\rho < 0$. Cela implique en effet d'inverser les amplitudes de l'impulsion initiale $\epsilon_1$ tout en les diminuant. Le phenomene se repete sur Fig. \ref{fig:r500}, bien que l'amplitude maximale semble attendre la tension d'entree en 10 $V_p_p$ a mesure que $Z_L$ tend vers l'infini (voir Tab. \ref{tab:amp}). 
\begin{figure}[h!]
    \centering
    % Première figure
    \begin{minipage}[b]{0.48\linewidth}
        \centering
        \includegraphics[width=\linewidth]{electrique/r100.png}
            \caption{$Z_L$ = 100 $\ohm$}
        \label{fig:r100}
    \end{minipage}
    \hfill
    % Deuxième figure
    \begin{minipage}[b]{0.48\linewidth}
        \centering
        \includegraphics[width=\linewidth]{electrique/r500.png}
            \caption{$Z_L$ = 500 $\ohm$}
        \label{fig:r500}
    \end{minipage}
\end{figure}

Les amplitudes maximales et minimales pour chaque intervale sont visible en annexe sur Tab. \ref{tab}. Nous pouvons determiner les coefficients de reflexion a partir de Eq. \ref{eq:}, soit par divisant l'amplitude de rebond de la source  par l'amplitude source ainsi que le coefficient de transmission en divisant l'amplitude de transmise par l'amplitude source. Or les coefficients de reflexion et de transmission de l'analyse sont pour l'instant effectuee a partir de la relation Eqs. \ref{eq:} et \ref{eq:}. 

\newpage
\subsection{Analyse}
Si l'on compare avec la theorie (voir Tab. \ref{tab:th}), nos valeurs pratiques sur Tab. \ref{eq:pr} sont ... :
\begin{table}[ht]
\centering
\begin{tabular}{|c|c|c|}
\hline
	\textbf{Lien $Z_C , Z_L$} & \textbf{$\rho_t_h$} & \textbf{$\rho_m_e_s$} \\
\hline
	$Z_L \approx Z_C$, $Z_L$ = 40 & $\rho \approx 0$ & -0.161  \\%Adaptation, pas de réflexion \\
        :wq
\hline
$Z_L < Z_C$, $Z_L$ = 25 & $\rho < 0$ & -0.377 \\%Réflexion avec inversion de phase \\
\hline
	$Z_L > Z_C$, $Z_L = 100$ & $\rho > 0$ & 0.288 \\%Réflexion sans inversion de phase \\
\hline
	$Z_L \rightarrow \infty$, $Z_L = 2500$ & $\rho \rightarrow +1$ & 0.957 \\
\hline
$Z_L \rightarrow 0$, $Z_L = 5$ & $\rho \rightarrow -1$ & -0.834 \\
\hline
\end{tabular}
	\label{tab:th}
\caption{Coefficients de reflexion et de transmission avec leur couple d'impedance}
\end{table}

Nous pouvons verifier l'evolution de $\rho$ en fonction de $Z_L$ sur Fig. \ref{fig:rhoo} trace a partir de la relation Eq. \ref{eq:rho}.

    \begin{figure}[ht]
        \centering
        \includegraphics[scale=0.38]{rho.png}
	    \caption{$\rho$($Z_L$)}
        \label{fig:rhoo}
    \end{figure}

\section{Conclusion}

Dans la partie acoustique, la mesure de la célérité des ultrasons dans l’air a conduit à une valeur $c = (3,3 \pm 0,4)\cdot 10^2,\mathrm{m/s}$, en bon accord avec l’ordre de grandeur attendu compte tenu des incertitudes expérimentales. L’étude de la dépendance angulaire et de la décroissance de l’amplitude en fonction de la distance a permis de confirmer la nature sphérique des ondes acoustiques, ainsi que la loi de décroissance de la pression en (1/r) avec des pentes lineaires de $\alpha_U = 155 \pm 8$ V/m, $\alpha_U_^2$ $= 30067 \pm 1771$ $V^2$/$m^2$, $\alpha_U_i_n_v = 0.21 \pm 0.02$ $V/m$ (voir Figs. \ref{fig:U(1/R)}, \ref{fig:U2(1/R2)}, \ref{fig:U(1/R)inv}).

Dans la partie électrique, l’étude de la propagation d’une onde de tension le long d’un câble coaxial a permis de déterminer le temps de propagation de $\tau \approx 0.5 \mu$s. A partir des mesures de la capacité et de l’inductance linéiques, nous avons determine l’impédance caractéristique réelle du câble  $Z_C = 55,5 \pm 0,3 \ohm$. La valeur d'impedance calculee est proche de la valeur nominale du câble RG58, ce qui confirme la pertinence de la méthode employée. Les coefficients de réflexion et de transmission n'ont pas été mesurés a partir des amplitudes, mais des impédances. Cela est dû au manque de temps.


